\documentclass[12pt,a4paper]{article}
\usepackage{amsmath,amsthm,amssymb}
\usepackage[utf8]{inputenc}
\usepackage[russian]{babel}

\makeatletter

\textheight=251mm %+20mm
\textwidth=163.4mm
\oddsidemargin=-2.1mm
\evensidemargin=-2.1mm
\topmargin=-16.5mm %-20mm

\newcommand*{\ralph}[1]{\expandafter\rralph\csname c@#1\endcsname}
\newcommand*{\rralph}[1]%
{\ifcase #1\or а\or б\or в\or г\or д\or е\or ж\or з\or и\or
к\or л\or м\or н\or о\or п\or р\or с\or т\or у\or
ф\or х\or ц\or ч\or ш\or щ\or э\or ю\or
я\fi}

\newcounter{problem}
\newcounter{item}
\newcommand*{\z}{\par\medskip\addtocounter{problem}{1}\setcounter{item}{0}
\textbf{Задача \arabic{problem}.\ }}
\newcommand*{\zh}{\par\medskip\addtocounter{problem}{1}\setcounter{item}{0}
\textbf{Задача \arabic{problem}*.\ }}
\newcommand*{\zH}{\par\medskip\addtocounter{problem}{1}\setcounter{item}{0}
\textbf{Задача \arabic{problem}**.\ }}
\newcommand*{\zn}[1]{\par\medskip\addtocounter{problem}{1}\setcounter{item}{0}
\textbf{Задача \arabic{problem} (#1).}}
\newcommand*{\znh}[1]{\par\medskip\addtocounter{problem}{1}\setcounter{item}{0}
\textbf{Задача \arabic{problem}* (#1).}}
\newcommand*{\znH}[1]{\par\medskip\addtocounter{problem}{1}\setcounter{item}{0}
\textbf{Задача \arabic{problem}** (#1).}}

\newcommand*{\zite}{\par\medskip\addtocounter{problem}{1}\setcounter{item}{1}
\textbf{Задача \arabic{problem}.} (\ralph{item})\ }
\newcommand*{\ite}{\par\addtocounter{item}{1}(\ralph{item})\ }
\newcommand*{\iteh}{\par\addtocounter{item}{1}(\ralph{item})*\ }
\newcommand*{\iteH}{\par\addtocounter{item}{1}(\ralph{item})**\ }

\newcommand*{\iten}[1]{\par\addtocounter{item}{1}(\ralph{item}) \textbf{(#1)}}
\newcommand*{\itenh}[1]{\par\addtocounter{item}{1}(\ralph{item})* \textbf{(#1)}}
\newcommand*{\itenH}[1]{\par\addtocounter{item}{1}(\ralph{item})** \textbf{(#1)}}
\newcommand{\myP}[0]{\mathsf{P}}

\begin{document}
\pagestyle{empty}

%\renewcommand{\@oddhead}%
%{\vbox{
{\footnotesize\begin{tabbing}V зимняя школа <<Комбинаторная математика и теория алгоритмов>>\`10--17 февраля 2013 г.\\
Виталий Павленко, ФИВТ МФТИ\`\texttt{vk.com/vitalypavlenko}\end{tabbing}}

%}}

\begin{center}
\vskip 24pt
{\bf\LargeТеорема Байеса -- 1: спам-фильтр}
\end{center}

Для получения зачёта по спецкурсу нужно решить все или почти все задачи из~двух листочков. Первый листочек посвящен наивному байесовскому классификатору, второй будет посвящен условной энтропии. Рассказывать решения задач можно мне в~любое время.

\z Имеются две монеты: честная и нечестная. У честной монеты вероятность выпадения решки равна $0.5$. У нечестной~--- $0.1$. Сначала мы равновероятно выбираем одну из двух монет, после чего кидаем её. Пусть при этом броске выпала решка. Какова вероятность того, что мы бросали честную монету?

\z Пусть $A$ и $B$~--- два события на некотором вероятностном пространстве~$\Omega$. Пусть $\myP(B) > 0$. Тогда события $A$ и $B$ независимы тогда и только тогда, когда ${\myP(A \mid B) = \myP(A)}$.

\z Бросают два игральных кубика: белый и красный. Посчитайте условную вероятность того, что на каждом кубике выпало не менее 3 очков, при условии, что сумма выпавших очков равна 7.

\z Ваша задача~--- построить классификатор, который будет отличать английские слова от не-английских. Слова порождаются источником со следующими вероятностями появления слов:

\vskip 4pt
\begin{tabular}{llll}
элемент в $\Omega$ & слово & английское слово? & вероятность \\
\hline
1 & ozb & нет & 2/7 \\ 
2 & uzu & нет & 2/7 \\
3 & zoo & да & 3/14 \\
4 & bun & да & 3/14 \\
\end{tabular}
\vskip 4pt

\ite Вычислите вероятности классов и условные вероятности для multinomial-\linebreak классификатора\footnote[1]{Multinomial-классификатор~--- это классификатор, который учитывает не только наличие/отсутствие признака, но и его количество. Например, если какая-то буква в слове встречается дважды, то это даст два множителя в итоговой формуле классификатора.}, который использует в качестве признаков буквы b, n, o, u и z. Исходите из того, что выборка, на которой будет обучаться классификатор, будет идеально отражать вероятности источника слов. Также примите в качестве предположения, что позиция буквы в слове не влияет на его английскость. Используйте сглаживание, при котором вычисленные нулевые вероятности заменяются вероятностью $0.01$, а ненулевые вероятности не изменяются. (При этом может оказаться, что $\myP(A) + \myP(\overline A) > 1$. С~математической точки зрения это абсурдно, но при построении классификатора на это можно закрыть глаза.)

\ite Считает ли классификатор слова \textsf{zoo}, \textsf{buzz}, \textsf{noob} английскими?

\ite Пусть теперь классификатор строится так, что признаком является не буква, а~пара (буква, позиция). К какому классу новый классификатор относит слово \textsf{zoo}? Вам нужно посчитать не все условные вероятности, а только те, которые влияют на~классификацию слова \textsf{zoo}.

\end{document}